% SHARD-ID: AISM-01-CLASSICAL-EQUIV
% SHARD-TITLE: The signed-stochastic bridge
% SHARD-SUMMARY: Reproduces lem-classical-equiv, the af-validated equivalence of the signed and stochastic formulations up to universal constants.
% SHARD-SUMMARY: Gives a prose account of the binomial-series construction forward and the positive-part renormalisation backward, with the explicit constants of the tree.
% SHARD-KEYWORDS: lem-classical-equiv, af validated, signed picture, stochastic picture, bridge, binomial series
\section{The signed-stochastic bridge}\label{sec:classical-equiv}

\begin{lemma}[\texttt{lem-classical-equiv}]\label{lem:classical-equiv}
The signed-idempotent and stochastic-idempotent formulations of classical stability are
equivalent up to universal constants. There is a universal $C>0$ such that:
\begin{enumerate}
\item (stochastic $\to$ signed) if $Q$ is row-stochastic with
  $\opnorm{Q^2-Q}\le\defect$, then $P=\theta(2Q-1)$ is a signed affine retraction
  (\texttt{def-signed-idempotent}) satisfying $\opnorm{P-Q}\le C\defect$ and negative
  mass $\negmass\le C\defect$;
\item (signed $\to$ stochastic) conversely, row-normalising the positive parts
  $\posp{p_i}$ of a signed idempotent $P$ yields a row-stochastic $Q$ with
  $\opnorm{P-Q}\le 2\negmass$ and $\opnorm{Q^2-Q}\le 6\negmass+4\negmass^2$.
\end{enumerate}
\emph{Transcription note.} The contract writes the map norm as $\lVert\cdot\rVert$
without a subscript; per \texttt{CONVENTIONS.md}~(b) it is the $\infty\!\to\!\infty$ map
norm $\opnorm{\cdot}$. The operator $\theta$ is carried verbatim from the contract,
which does not expand it; the validated tree instantiates it as the functional calculus
$\theta(X)=\tfrac12(I+\sgn X)$ of \texttt{def-theta-idempotent-approximation}.
\end{lemma}

\noindent\textbf{Registry contract (verbatim).}
\contractquote{The signed-idempotent and stochastic-idempotent formulations of classical stability are equivalent up to universal constants: Q row-stochastic with ||Q^2-Q|| <= eta gives P=theta(2Q-1) signed affine retraction with ||P-Q|| <= C eta and neg mass delta <= C eta, and conversely row-normalising p_i^+ gives Q with ||P-Q|| <= 2 delta, ||Q^2-Q|| <= 6 delta+4 delta^2.}

\paragraph{Proof (prose account of the af-validated tree).}
\emph{Forward direction.} Put $X=2Q-I$. Row-stochasticity gives $X\mathbf1=\mathbf1$ and
$\opnorm{X}\le3$, while the algebraic identity
\begin{equation}\label{eq:ce-square}
  X^2-I=4\,(Q^2-Q)
\end{equation}
converts the almost-idempotence hypothesis into $\opnorm{X^2-I}\le4\defect$. Write
$Y=X^2-I$. Two properties of $Y$ drive everything: $Y\mathbf1=0$ (because
$X\mathbf1=\mathbf1$ forces $X^2\mathbf1=\mathbf1$), and $\opnorm{Y}\le4\defect\le\tfrac12$
once $\defect\le\tfrac18$.

The smallness of $Y$ puts the inverse square root inside its convergence radius: the
binomial series
\[
  B=(I+Y)^{-1/2}=\sum_{m\ge0}\binom{-\tfrac12}{m}Y^{m}
\]
converges absolutely in the map norm. Being a norm limit of polynomials in $Y=X^2-I$, the
operator $B$ commutes with $X$; since $Y\mathbf1=0$ kills every positive power,
$B\mathbf1=\mathbf1$; and the absolute coefficient sum gives the quantitative bound
$\opnorm{B-I}\le(1-t)^{-1/2}-1\le 2t$ at $t=\opnorm{Y}\le\tfrac12$, hence
$\opnorm{B-I}\le8\defect$. The scalar identity $\bigl((1+z)^{-1/2}\bigr)^2(1+z)=1$
transfers to the absolutely convergent series in $Y$, so $B^2X^2=I$; commutation then
makes $S=XB$ an exact involution,
\[
  S^2=XBXB=X^2B^2=I,\qquad S\mathbf1=XB\mathbf1=X\mathbf1=\mathbf1,
\]
with $S-X=X(B-I)$ and therefore $\opnorm{S-X}\le3\cdot8\defect=24\defect$.

Setting $P=\theta(X)=\tfrac12(I+S)$ now gives $P\mathbf1=\mathbf1$ and
$P^2=\tfrac14(I+2S+S^2)=\tfrac12(I+S)=P$, so $P$ is an exact signed idempotent in the
sense of \texttt{def-signed-idempotent}. Since $Q=\tfrac12(I+X)$, the difference is
$P-Q=\tfrac12(S-X)$ and hence $\opnorm{P-Q}\le12\defect$. Finally the negative mass:
row-stochastic $Q$ is entrywise nonnegative (\texttt{def-stochastic}), so for every entry
$\max\{-P_{ij},0\}\le\lvert P_{ij}-Q_{ij}\rvert$ --- trivially where $P_{ij}\ge0$, and
because $-P_{ij}\le Q_{ij}-P_{ij}$ where $P_{ij}<0$. Summing over $j$ and maximising over
$i$ gives, with $\negmass$ as in \texttt{def-negative-mass},
\begin{equation}\label{eq:ce-negmass}
  \negmass(P)\;\le\;\opnorm{P-Q}\;\le\;12\defect .
\end{equation}
So $C=12$ serves. The export records an instructive scope correction here: the
negative-mass step was first blocked because \texttt{def-negative-mass} --- although
permitted by the registry shard --- had not been registered inside the workspace, and the
tree carried an explicit gap certificate rather than the estimate; registering the
definition discharged it.

\emph{Converse direction.} Let $P$ be an exact signed idempotent with negative mass
$\negmass$. Split each row into disjointly supported nonnegative parts,
$p_i=\posp{p_i}-\negp{p_i}$, and set $a_i=\sum_k(P_{ik})_-\le\negmass$. Total mass one
(\texttt{def-signed-idempotent}) gives $\sum_k\posp{p_{ik}}=1+a_i$, so
\[
  q_i=\frac{\posp{p_i}}{1+a_i}
\]
is a probability vector and the matrix $Q$ with rows $q_i$ is row-stochastic. The
correction $e_i=q_i-p_i$ is
\[
  e_i=\negp{p_i}-\frac{a_i}{1+a_i}\,\posp{p_i},
\]
a difference of two disjointly supported vectors of $\ell^1$-mass $a_i$ each; hence
$\lone{e_i}=2a_i\le2\negmass$ and $\opnorm{P-Q}\le2\negmass$.

For the idempotence defect, write $D=Q-P$. The row-geometry clause of
\texttt{def-signed-idempotent} gives $\opnorm{P}\le1+2\negmass$, hence
$\opnorm{P-I}\le2+2\negmass$. Using $P^2=P$ exactly, each row satisfies the exact
decomposition
\begin{equation}\label{eq:ce-defect}
  q_iQ-q_i=e_i(P-I)+q_iD .
\end{equation}
Since $q_i$ is a probability vector, $\lone{q_iD}\le\opnorm{D}\le2\negmass$, while
$\lone{e_i(P-I)}\le2\negmass\,(2+2\negmass)$. Adding,
$\opnorm{Q^2-Q}\le6\negmass+4\negmass^2$, exactly as contracted.

\medskip
\noindent\textbf{Status.} \texttt{proved}; \texttt{af: validated} (T0). Workspace
\texttt{proofs/lem-classical-equiv}: a 29-node tree (21 live nodes validated, 8
superseded nodes archived), root \texttt{validated}, taint clean.

\noindent\textbf{Role in the chain.} This is the licence to move between the two
pictures of \S\ref{sec:overview}, and the only such licence in the registry. It has no
dependencies. Everything stated in the signed picture reaches the stochastic
formulation of \texttt{op-classical} only through it, and only up to the universal
constants above; \texttt{op-classical} itself is now discharged at the af-validated rung via Route F (\S\ref{sec:status-outlook}).
